\section*{摘\quad 要}
\addcontentsline{toc}{section}{摘要}

携带卫星遥感大数据先验底图的自主水下航行器（AUV）在动态海洋环境中执行应急巡测时，面临机载计算能力、电池能量储备与作业时间窗口的多重物理约束。在先验底图因云层遮挡与海流平流而不可避免发生时效性退化的条件下，航行器必须在每个决策周期内反复权衡时间与能量的分配：是唤醒高功耗机载图形处理器（GPU）在内部底图上深入思考规划，还是维持低功耗巡航开展现场采样，抑或关停推进器随流漂移以开展原位重采样。既有移动观测方法多依赖经验设定的固定资源分配规则，难以适应不同机载算力速度与环境底图误差水平；而经典单步元推理方法直接利用规划器内部搜索轨迹评估计算价值，在有限样本集合上易受优化器诅咒诱导而陷入过度思考，并忽视了全任务周期内能量消耗与空间探索的非线性耦合。针对上述局限，本文提出一种基于留出假想海洋的集合同化与行前预演元控制框架，在任务出发前通过完整闭环假想任务推演实现思考与行动资源的自适应分配。该方法首先利用多年历史卫星遥感大数据校准不同滞后天数的底图误差统计特征，围绕滞后底图生成48个空间相关的假想海洋样本作为先验集合，并在随流漂移参考系下采用带局地化函数的序列集合调整卡尔曼滤波同化现场荧光计观测；其次，将假想海洋严格解耦为32个供交叉熵规划器搜索的样本与16个从未参与搜索的留出验证样本，以消除轨迹评估中的过拟合偏差；最后，在航行器出发前抽取4个留出假想海洋作为合成真实海洋，在匹配真实机载算力速度与电池功耗模型的条件下，采用共同随机数并行预演五种典型时间分配策略的24小时完整任务过程，择优执行平均随时遗憾值最低的策略。本文基于2018至2026年东海海域3197天多源融合叶绿素卫星遥感大数据，构建422个春季有害赤潮应急巡测任务窗口（含2018至2022年274个开发窗口与2023至2026年148个测试窗口），对所提方法开展系统实验验证。实验结果表明：当单步规划预算从64次推演扩展至200,000次时，真实海洋路线收益在约2,000次推演后趋于饱和（从0.749增至0.793并收敛于0.803），规划器内部搜索得分却持续攀升至0.917（虚高0.114），而16个留出假想海洋的评估均值在全预算跨度内与真实海洋收益之差始终保持在0.005以内；高功耗深度思考仅在机载基础算力相对规划问题过慢时产生显著正收益，巡航途中停航漂移重采样在所有测试工况下均呈负收益；在148个测试任务窗口上，经典单步元推理因65\%周期启动高功耗思考且28\%周期停航思考，在慢速与中速算力平台上的随时遗憾值较特定硬件事后最优调参策略分别劣化0.017与0.021，而本文提出的行前假想预演在慢速与中速算力平台上分别取得0.276与0.265的随时遗憾值，与特定硬件最优调参策略持平（差值分别为$+0.007$与$-0.001$）。本研究从机理上揭示了具身智能体轨迹规划中的过思考偏差及其集合解耦校正律，为海洋大数据赋能的实时科学计算与自主移动观测系统提供了无需历史任务调参的自校准元控制范式。

\vspace{1em}
\noindent\textbf{关键词：}海洋大数据；自主水下航行器；元推理；数据同化；赤潮巡测；计算与行动权衡

\clearpage
\section*{Abstract}
\addcontentsline{toc}{section}{Abstract}

Autonomous underwater vehicles (AUVs) equipped with prior maps derived from satellite ocean big data operate under tight physical constraints of onboard compute, battery capacity, and mission time when conducting emergency marine surveys. Because satellite prior maps inevitably degrade due to cloud gaps, retrieval uncertainties, and advective evolution, the vehicle must repeatedly decide at every control interval how to allocate its finite time and energy: whether to activate a power-hungry onboard graphics processing unit (GPU) to think more deeply on its internal map, continue cruising and sampling under low-power background planning, or shut down propulsion to drift and resample locally. Existing adaptive sampling approaches typically rely on fixed allocation heuristics that fail to generalise across varying onboard processor speeds and map error regimes, whereas classic one-step metareasoning estimates the value of computation from the planner's own internal improvement trace, succumbing to the optimizer's curse and chronic overthinking while ignoring full-mission coupling between battery depletion and spatial coverage. To overcome these limitations, this paper proposes an ensemble data assimilation and pre-mission rehearsal meta-control framework based on held-out imagined oceans, enabling self-calibrated allocation of thinking and acting resources prior to deployment. Specifically, historical satellite archives are used to calibrate stale-map error statistics and generate an ensemble of 48 spatially correlated imagined oceans around the stale satellite map, which are sequentially updated in a water-drifting reference frame using a localised Ensemble Adjustment Kalman Filter as in-situ fluorometer observations arrive; the ensemble is strictly split into 32 searched members used by a warm-started cross-entropy trajectory planner and 16 held-out members reserved exclusively for unbiased value estimation; before setting off, the vehicle treats 4 held-out imagined oceans as synthetic truths and rehearses complete 24-hour missions across five candidate time-allocation strategies under its exact hardware speed and battery model with common random numbers, executing the strategy that minimises imagined time-averaged regret. We evaluate the framework using 3,197 days of gap-filled merged satellite chlorophyll-a data over the East China Sea from 2018 to 2026 across 422 spring harmful algal bloom survey windows (274 development windows in 2018--2022 and 148 test windows in 2023--2026). Results show that as the planning budget scales from 64 to 200,000 rollouts per decision, true route value in the real ocean saturates near 2,000 rollouts (rising from 0.749 to 0.793 and plateauing at 0.803) while the planner's internal search score keeps climbing to 0.917 (overestimating reality by 0.114), whereas evaluation on the 16 held-out imagined oceans remains within 0.005 of reality across all budgets; heavy GPU thinking pays off only when the baseline processor is too slow for the planning horizon, and stopping mid-mission to resample never pays in this domain; on the 148 test windows, classic one-step metareasoning spends 65\% of intervals thinking on the GPU and 28\% stopped to think, degrading time-averaged regret by 0.017 to 0.021 relative to per-vehicle tuned strategies, whereas pre-mission rehearsal achieves time-averaged regrets of 0.276 and 0.265 on slow and moderate compute platforms respectively, matching oracle per-vehicle tuning within $+0.007$ and $-0.001$. This work elucidates the mechanism of overthinking bias in embodied trajectory planning and establishes a tuning-free, self-calibrating meta-control paradigm for ocean-big-data-driven autonomous observation.

\vspace{1em}
\noindent\textbf{Keywords: }Ocean big data; Autonomous underwater vehicle; Metareasoning; Data assimilation; Harmful algal bloom survey; Thinking-acting trade-off
